\documentclass[12pt]{article}
\usepackage[utf8]{inputenc}
\usepackage[T1]{fontenc}
\usepackage{hyperref}
\title{Muluridyi Language Resources for K-12 Curriculum Development}
\author{Compiled by Claude Sonnet 5 (claude-sonnet-5) for Andr\'as Kornai}
\date{}
\begin{document}
\maketitle

\section{Introduction}
Muluridyi (ISO 639-3 vmu; Glottolog mulu1243) is a Pama-Nyungan language of north Queensland, Australia. It is classed as extinct. Searches found no teaching materials, word lists or revitalization programme specific to it. This dossier is therefore very short.

\section{Primary Educational Resources}
None were found in search results.

\section{Linguistic and Reference Materials}
\begin{itemize}
\item AIATSIS Austlang entry for Muluridji. It notes that the variety is treated as a dialect of Kuku Yalanji, so Kuku Yalanji documentation may hold related material: \url{https://collection.aiatsis.gov.au/austlang/language/y97}
\item South Australian Museum archive entry on the Muluridji language group, with its traditional territory: \url{https://www.samuseum.sa.gov.au/collection/archives/language_groups/muluridji}
\item Glottolog classification (Pama-Nyungan, Yalandyic): \url{https://glottolog.org/resource/languoid/id/mulu1243.bigmap.html}
\end{itemize}

\section{Religious and Community Resources}
None were found.

\section{Digital and Audio Resources}
\begin{itemize}
\item Endangered Languages Project entry, with metadata only: \url{https://endangeredlanguages.com/lang/3373}
\item Wikitongues language page, with metadata only: \url{https://wikitongues.org/languages/vmu/}
\end{itemize}

\section{Curriculum Development Context}
No Muluridyi lexicon or recordings were found in the searches. A heritage curriculum would have to start from archival work, for example by contacting AIATSIS and by consulting the Kuku Yalanji literature (Patz). It should also involve local Aboriginal communities. This was not verified here.
\end{document}